\subsection{Kernel output package}

The benchmark row for $\alpha^{-1}(0)$ is defined by the positive self-consistent root of \Cref{eq:carrier-cfe}, not by freezing the opening at its retained seed value. On the canonical branch that positive root remains
\[
\alpha^{-1}(0)=\num{137.03599921615843}.
\]
For the remaining one-seed relations we keep the retained branch opening
\begin{equation}
\label{eq:hard-kernel-values}
\cthree=\frac{1}{8\pi},
\qquad
\phiz=\frac{1}{6\pi}+\frac{3}{256\pi^4},
\end{equation}
\begin{remark}[Two retained openings and their roles]
The two symbols belong to the same retained branch but play different jobs. The electromagnetic $\alpha$ row uses the self-consistent retained opening $\phiztwo(\alpha)$ inside the carrier-form closure equation, whereas the one-seed bridge relations freeze that branch to the retained canonical seed $\phiz$. Thus $\phiztwo(\alpha)$ belongs only to the $\alpha$ benchmark line, while $\phiz$ is the seed used for $(\beta_{\mathrm{rad}},\Omega_b,\lambda_C,\sin^2\theta_{13})$.
\end{remark}
and the same retained kernel package emits the compact seed chain
\begin{equation}
\label{eq:kernel-benchmark-formulas}
C_{a\gamma\gamma}=-4\cthree=-\frac{1}{2\pi},
\qquad
\beta_{\mathrm{rad}}=\frac{\phiz}{4\pi},
\qquad
\Omega_b=(4\pi-1)\beta_{\mathrm{rad}},
\qquad
\lambda_C=\sqrt{\phiz(1-\phiz)},
\qquad
\sin^2\theta_{13}=\phiz e^{-\gamma}.
\end{equation}
\begin{remark}[Why the $\alpha$ row must stay self-consistent]
The benchmark number $\alpha^{-1}(0)=\num{137.03599921615843}$ comes from the retained second-order self-consistent equation with $\phiztwo(\alpha)$ and the term $\delta_2$. If one instead freezes the opening at
\[
\phiz=\frac{1}{6\pi}+\frac{3}{256\pi^4},
\]
the root shifts to $\alpha^{-1}(0)\approx \num{137.036501465}$, i.e. by about $\num{5.02e-4}$, which is roughly $\num{2.4e4}\sigma$ on the CODATA/NIST uncertainty. Even dropping only $\delta_2$ still shifts the benchmark by about $243\sigma$. The wording of the $\alpha$ benchmark must therefore remain explicit: full self-consistency and second order are load-bearing, not editorial garnish.
\end{remark}
\begin{tcolorbox}[colback=blue!3, colframe=blue!50!black, breakable,
    title={\textbf{Kernel-output package}}, fonttitle=\bfseries]
The point of this section is not to claim more numerics than TFPT currently owns. It is to show that the established kernel already fixes a coherent electromagnetic package before the later completion theorems begin: the seam constant $\cthree$, the retained canonical seed $\phiz$, the self-consistent carrier-form CFE for $\alpha$, the dimensionless axion--photon anomaly coefficient, the birefringence seed, and the Cabibbo anchor.
\end{tcolorbox}

\section{Compression identities and carrier exclusion}

This section collects the algebraic compressions that turn the carrier packet and kernel numerics from a list of internally consistent outputs into an overdetermined structure with genuine exclusion power. The compression identities are not new theorems; they are exact algebraic consequences of the already established carrier data. The exclusion theorems are the sharpest available evidence that the carrier packet actively selects the physical world.

\subsection{Rank lock and the \texorpdfstring{$2$-$3$-$5$}{2-3-5} arithmetic}

\begin{proposition}[Carrier rank lock]
\label{prop:rank-lock}
If the carrier identity $\dim S^+=2^{g-1}$, the canonical family data $N_{\mathrm{fam}}=3$, and the optional occupancy compression $\Oadm=(g-1)\dim\mathfrak g_{\SM}=12(g-1)$ are combined, then
\begin{equation}
\label{eq:rank-lock}
3\cdot 2^{g-1}=12(g-1)
\qquad\Longleftrightarrow\qquad
2^{g-3}=g-1.
\end{equation}
The unique integral solution is $g=5$.
\end{proposition}

\begin{proof}
Set $f(g):=2^{g-3}-(g-1)$. Then
\[
f(1)=\frac14,\qquad
f(2)=-\frac12,\qquad
f(3)=-1,\qquad
f(4)=-1,\qquad
f(5)=0.
\]
For $g\ge 4$,
\[
f(g+1)-f(g)=2^{g-3}-1>0,
\]
so $f$ is strictly increasing on $g\ge 4$. Since $f(4)<0$ and $f(5)=0$, there is no solution with $g\ge 6$. A direct check for $g\le 4$ shows no other solution. Hence the unique integral solution is $g=5$.
\end{proof}

\begin{remark}[The $2$-$3$-$5$ arithmetic]
The hard carrier packet rests on exactly three discrete numbers: $2$ (the double cover), $3$ (the color rank and family count), and $5$ (the carrier rank). From these three atoms one recovers
\[
\gamma=\frac{5}{6}=\frac{5}{2\cdot 3},
\qquad
B=\frac{3}{2},
\qquad
B\gamma=\frac{5}{4}=\frac{5}{2^2},
\]
\[
\Oadm=3\cdot 2^4=48,
\qquad
\Dstart=5\cdot 12=60,
\qquad
\frac{\Dstart}{\Oadm}=\frac{5}{4}=B\gamma,
\]
and the defect coefficients $\deltatop=48\,\cthree^4$ and $\delta_2=\frac54\,\deltatop^2=2880\,\cthree^8$ with $2880=60\cdot 48$.
The persistent appearance of $5/4$ as carrier asymmetry $\times$ norm, second defect stage, and ladder-to-occupancy ratio is therefore not three coincidences but one rational shadow of the $2$-$3$-$5$ triad.
\end{remark}

\begin{lemma}[Universal \texorpdfstring{$5/4$}{5/4} compression quotient on the minimal branch]
\label{lem:universal-5-4-compression}
Define
\[
\rho_{5/4}:=B\gamma.
\]
On the canonical minimal branch,
\[
\rho_{5/4}
=
B\gamma
=
\frac{\delta_2}{\deltatop^2}
=
\frac{\Dstart}{\Oadm}
=
\frac{\gcar}{\gcar-1}
=
\frac54.
\]
Thus $5/4$ is the universal compression quotient of the minimal branch rather than a recurring numerical motif.
\end{lemma}

\begin{proof}
The identity $B\gamma=\frac54$ is \Cref{cor:minimal-carrier-compression}. The defect relation gives
\[
\frac{\delta_2}{\deltatop^2}=\frac54.
\]
On the canonical branch,
\[
\frac{\Dstart}{\Oadm}=\frac{60}{48}=\frac54,
\]
and with $\gcar=5$ one also has
\[
\frac{\gcar}{\gcar-1}=\frac{5}{4}.
\]
Hence all displayed ratios agree on the same exact value.
\end{proof}

\begin{remark}[Interpretive status of the quotient]
The quotient itself is exact. Any holographic or spacetime-dimensional reading of the denominator $\gcar-1=4$ remains interpretive in the present manuscript rather than theorem-level.
\end{remark}

\begin{remark}[Conditional rank-5 master compression]
If one accepts the carrier rank as the master parameter, then on the canonical branch
\[
N_{\mathrm{fam}}=\frac{\gcar+1}{2}=3,
\qquad
\Oadm=N_{\mathrm{fam}}\,2^{\gcar-1}=\frac{(\gcar+1)}{2}\,2^{\gcar-1}=(\gcar+1)\,2^{\gcar-2}=48,
\]
\[
\Oadm\gamma=\gcar\,2^{\gcar-2}=40,
\qquad
\bone=\frac{\gcar\,2^{\gcar-2}+1}{10}=\frac{41}{10}.
\]
This is not yet a theorem. But if $N_{\mathrm{fam}}=(\gcar+1)/2$ closes geometrically, then family count, occupancy, abelian index, and the $41$-chain all collapse onto the single integer $\gcar=5$.
\end{remark}

\begin{corollary}[Pythagorean compression of the integrity index on the rank-locked branch]
\label{cor:pythagorean-integrity-index}
Assume the conditional compression branch
\[
10\,\bone=\gcar\,2^{\gcar-2}+1
\]
and the rank lock
\[
2^{\gcar-3}=\gcar-1.
\]
Then
\[
\mathcal I_{41}
=
10\,\bone
=
2\gcar(\gcar-1)+1
=
\gcar^2+(\gcar-1)^2.
\]
In particular, on the minimal rank-locked branch $\gcar=5$ one gets
\[
\mathcal I_{41}=5^2+4^2=41.
\]
\end{corollary}

\begin{proof}
Multiplying the rank lock by $2\gcar$ gives
\[
\gcar\,2^{\gcar-2}=2\gcar(\gcar-1).
\]
Inserting this into the conditional compression formula yields
\[
10\,\bone=2\gcar(\gcar-1)+1.
\]
Since
\[
2\gcar(\gcar-1)+1=\gcar^2+(\gcar-1)^2,
\]
the claim follows.
\end{proof}

\begin{remark}[The number $4$ as a cross-cutting echo]
In addition to the $2$-$3$-$5$ triad, the integer $4$ appears at four structurally distinct levels:
\[
F_{\mathrm{patch}}=4,
\qquad
\gcar-1=4,
\qquad
\frac{\Oadm}{\dim\mathfrak g_{\SM}}=\frac{48}{12}=4,
\qquad
\frac{|R(E_8)|}{\Dstart}=\frac{240}{60}=4.
\]
That is, four corner patches, carrier rank minus one, matter-to-gauge compression factor, and root-to-start quotient all evaluate to the same integer.
\end{remark}

\subsection{Exact algebraic compressions on the minimal branch}

\begin{corollary}[Two-point spectral compression of $\gamma$ and $\phibase$]
Let
\[
y_+:=\frac12,
\qquad
y_-:=-\frac13
\]
be the two distinct carrier eigenvalues singled out by the factorized carrier polynomial. Then
\[
\gamma=y_+-y_-,
\qquad
\phibase=\frac{y_++y_-}{\pi}.
\]
Here $y_++y_-$ is the unweighted Vieta sum of the two distinct carrier eigenvalues, not the trace on $E$.
\end{corollary}

\begin{proof}
By \Cref{eq:hard-kernel-values} and the factorized carrier polynomial, the two distinct spectral values are exactly $y_+=1/2$ and $y_-=-1/3$. Their difference is
\[
y_+-y_-=\frac12+\frac13=\frac56=\gamma,
\]
while their sum is
\[
y_++y_-=\frac12-\frac13=\frac16.
\]
Since $\phibase=(1-\gamma)/\pi=1/(6\pi)$, the claim follows.
\end{proof}

\begin{corollary}[Family hypercharge variance and electroweak block trace]
Equip one carrier family with the uniform average
\[
\mathbb E_{S^+}(A):=\frac{1}{16}\tr_{S^+}A.
\]
Then
\[
\operatorname{Var}_{S^+}(Y)
:=
\mathbb E_{S^+}(Y^2)-\mathbb E_{S^+}(Y)^2
=
\frac{1}{16}\tr_{S^+}Y^2
=
\frac{5}{24}.
\]
Consequently,
\[
I_1^{\mathrm{EW}}=\operatorname{Var}_{S^+}(Y)\,\bone.
\]
\end{corollary}

\begin{proof}
The one-family trace identities give $\tr_{S^+}Y=0$ and $\tr_{S^+}Y^2=10/3$. Hence the centered second moment is
\[
\frac{1}{16}\tr_{S^+}Y^2=\frac{10/3}{16}=\frac{5}{24}.
\]
Using $I_1^{\mathrm{EW}}=\frac{5}{24}\,\bone$ from the electroweak block relation yields the second identity.
\end{proof}

\begin{corollary}[Factorial quadratic trace and two-defect factorization]
\label{cor:factorial-trace-delta2}
Let
\[
X:=6Y.
\]
Then the one-family quadratic trace is
\[
\tr_{S^+}(X^2)
=
36\,\tr_{S^+}(Y^2)
=
36\cdot\frac{10}{3}
=
120
=
5!.
\]
Equivalently, using the six hypercharge sectors with multiplicities
\[
(1,3,6,1,2,3)
\]
at
\[
X\in\{0,-4,1,6,-3,2\},
\]
the same trace is
\[
0^2+3\cdot 4^2+6\cdot 1^2+1\cdot 6^2+2\cdot 3^2+3\cdot 2^2=120.
\]
On the retained branch this gives the exact factorization
\[
\delta_2
=
2880\,\cthree^8
=
4!\,\tr_{S^+}(X^2)\,\cthree^8.
\]
\end{corollary}

\begin{proof}
By \Cref{prop:split-weight-enumerator}, one has
\[
\tr_{S^+}(Y^2)=\frac{10}{3}.
\]
Since $X=6Y$, the first identity follows immediately. The sector-by-sector sum is the same quadratic trace written in the basis of admissible occupancy sectors from \Cref{prop:family-hypercharge-compression}. Finally, \Cref{eq:hard-topological-kernel} gives
\[
\delta_2=2880\,\cthree^8=24\cdot 120\,\cthree^8
=
4!\,\tr_{S^+}(X^2)\,\cthree^8.
\]
\end{proof}

\begin{remark}[Quadratic trace versus variance]
The value $120=5!$ is the unnormalized quadratic trace $\tr_{S^+}(X^2)$, not the normalized variance. The corresponding uniform variance is
\[
\operatorname{Var}_{S^+}(X)=36\,\operatorname{Var}_{S^+}(Y)=36\cdot\frac{5}{24}=\frac{15}{2}.
\]
\end{remark}

\begin{remark}[Arithmetic Higgs lock on the canonical family branch]
Define the generic occupancy
\[
\Omega_{\mathrm{occ}}:=16\,N_{\mathrm{fam}}.
\]
Then the abelian-index formula gives
\[
10\,\bone
=
\frac{40}{3}N_{\mathrm{fam}}+N_\Phi
=
\Omega_{\mathrm{occ}}\,\gamma+N_\Phi.
\]
On the canonical family branch $N_{\mathrm{fam}}=3$, hence $\Omega_{\mathrm{occ}}=48$ and therefore
\[
10\,\bone=40+N_\Phi.
\]
Thus the $41$-chain arithmetically selects $N_\Phi=1$. This is an arithmetic lock on the canonical branch, not the dynamical Higgs-selection theorem, which in the present manuscript still comes from the bosonic relative-index / Callias route.
\end{remark}

